% Generated draft: requires mathematical and source review.
\documentclass[11pt]{article}
\usepackage[margin=25mm]{geometry}
\usepackage{fontspec}
\setmainfont{Latin Modern Roman}
\usepackage{amsmath,amssymb,mathtools,unicode-math}
\setmathfont{Latin Modern Math}
\usepackage{enumitem,needspace,xurl,hyperref}
\setlength{\parindent}{0pt}
\setlength{\parskip}{5pt}
\setlength{\emergencystretch}{3em}
\newcommand{\N}{\mathbb{N}}
\newcommand{\Z}{\mathbb{Z}}
\newcommand{\Q}{\mathbb{Q}}
\newcommand{\R}{\mathbb{R}}
\newcommand{\C}{\mathbb{C}}
\newcommand{\F}{\mathbb{F}}
\newcommand{\A}{\mathbb{A}}
\newcommand{\PP}{\mathbb{P}}
\newcommand{\E}{\mathbb{E}}
\newcommand{\eps}{\varepsilon}
\newcommand{\dd}{\,\mathrm{d}}
\begin{document}
\section*{UnsolvedMath ID: 800402}
\textbf{Source problem code:} catalogue entry
\subsection{Definitions and mathematical statement}
Let \(f:\mathbb{P}^{N}\dashrightarrow\mathbb{P}^{N}\) be a dominant rational map, meaning that its image is Zariski dense. Write \(f^{n}=f\circ\cdots\circ f\) for its \(n\)-fold iterate. For positive-dimensional projective space over a field, the algebraic degree \(\deg(f^{n})\) is the common degree of homogeneous coordinate polynomials representing \(f^{n}\), after their common factors have been removed. The supplied statement does not specify the base field or the permitted range of \(N\); these parameters require clarification.

A constant-coefficient linear recurrence sequence \((L_n)_{n\geq 1}\) is understood here to satisfy, for some fixed integer \(d\geq 1\) and fixed scalars \(c_0,\ldots,c_{d-1}\),
\[
L_{n+d}=\sum_{k=0}^{d-1}c_kL_{n+k}\qquad(n\geq 1).
\]
This uses the homogeneous interpretation of a linear recurrence. The source does not specify the scalar domain of the coefficients or sequence values, so no such domain is selected in this draft.

\textbf{Question (supplied source, Question 25).} For every dominant rational map \(f\) as above, do there exist an integer \(r\geq 1\), constant-coefficient linear recurrence sequences
\[
(L_n^{(1)})_{n\geq 1},\ldots,(L_n^{(r)})_{n\geq 1},
\]
and a function
\[
\iota:\{1,2,3,\ldots\}\longrightarrow\{1,\ldots,r\}
\]
such that
\[
\deg(f^{n})=L_n^{(\iota(n))}\qquad\text{for every }n\geq 1?
\]
The finite family and the index function may depend on \(f\). Different sequences may have different recurrence orders and coefficients. The source imposes no periodicity or other regularity condition on \(\iota\). The unspecified base field, dimension range, and recurrence scalar domain remain qualifications to this formulation.
\subsection{Short English statement}
Must the degrees of the iterates of every dominant rational self-map of projective space be obtainable by choosing, at each positive index, a term from one fixed finite family of constant-coefficient linear recurrence sequences? The source leaves the base field, dimension range, and recurrence scalar domain unspecified.
\subsection{Sources}
\begin{itemize}
\item Supplied local UnsolvedMath registry record: permanent ID \(800402\), exact source problem code \texttt{catalogue entry}, title ``Open Problems in Arithmetic Dynamics --- Question 25.'' The registry attributes the problem to Benedetto, DeMarco, Ingram, Jones et al., with proposed year 2018.
\item Benedetto, DeMarco, Ingram, Jones et al., \emph{Current Trends and Open Problems in Arithmetic Dynamics} (2018), Question 25. Supplied source URL: \url{https://arxiv.org/abs/1806.04980}. The supplied provenance records access on 2026-10-16 and extraction method \texttt{arxiv\_latex}; these are source-record claims, not a claim of independent access or verification.
\item The supplied registry labels the problem ``open'' and records an update on 2026-10-16. This is supplied status information, not an independently verified assessment of the literature.
\end{itemize}
\end{document}
